\appendix

% Séparation blanche entre l'exposé et les slides de réserve.
\begin{frame}[plain,noframenumbering]
\end{frame}

% RÉSERVE. Étude collective MoNNA, hors du périmètre personnel d'Arthur.
\begin{frame}{\monna{} décentralisé : étude de cas ICML 2023}
    \small
    L'équipe a réalisé cette simulation et analysé ses résultats.
    Mon travail portait sur le cas centralisé.

    MNIST réduit : $n=16$, $f=5$, 200 tours, \signflip{} d'amplitude 1.

    {\centering
        \includegraphics[width=\textwidth]{"MoNNA ICML 2023.png"}
        \par}

    La règle \nna{} moyenne les modèles les plus proches du modèle local.\\
    Exactitude finale sous attaque : \nna{} $\approx84\,\%$,
    \average{} $\approx10\,\%$. Sans attaque : $\approx87\,\%$.

    Farhadkhani et al., ICML 2023.
\end{frame}

% RÉSERVE. Extraits du listing de l'annexe 08 du mémoire.
\begin{frame}[fragile]{Une expérience centralisée}
    \small
    \begin{verbatim}
sim = KrumSimulation(
    model_cls=Krum2017MLPMnist,
    train_datasets=worker_datasets, test_set=test_set,
    aggregator=aggregator, attack=attack,
    n=n, f=f, rounds=100, batch_size=32, lr=0.1, seed=seed,
)
sim.setup()

loss = Metric("loss", dtype=float)
accuracy = Metric("accuracy", dtype=float)
for step in range(100):
    sim.step()
    if step % 10 == 0:
        test_loss, test_accuracy = sim.evaluate()
        loss.push(step, test_loss, skip_if_exists=True)
        accuracy.push(step, test_accuracy, skip_if_exists=True)
    \end{verbatim}

    La fonction d'évaluation calcule la perte et l'exactitude,
    c'est-à-dire la proportion de prédictions correctes.

    \medskip
    {\footnotesize Extrait du listing de l'annexe 08.}
\end{frame}

% RÉSERVE. Balayage de configurations tiré du listing de l'annexe 08.
\begin{frame}[fragile]{Faire varier les configurations d'une campagne}
    \small
    \begin{verbatim}
with Orchestrator("krum_byzantine_study") as orch:
    for n, f in [(13, 2), (25, 5)]:
        for aggregator in [Average, Krum, Bulyan]:
            for attack in [ALIEAttack, SignFlipAttack, None]:
                orch.run(
                    experiment, n=n, f=f,
                    aggregator=aggregator, attack=attack, seed=42,
                )

loss_df = orch.get("loss").to_pandas()
acc_df = orch.get("accuracy").to_pandas()
    \end{verbatim}

    \code{None} signifie qu'aucune attaque n'est appliquée.
    Les mesures sont converties en tableaux avec la bibliothèque pandas.
    L'orchestrateur a été développé par Sébastien Rouault et Peva Blanchard.

    \medskip
    {\footnotesize Extrait du listing de l'annexe 08.}
\end{frame}

% RÉSERVE. Notations des onze règles d'agrégation implémentées.
\begin{frame}{Notations des agrégateurs : moyenne et coordonnées}
    \small
    Soient $\vect{x}_1,\ldots,\vect{x}_n\in\mathbb{R}^d$ les gradients reçus.
    Pour une coordonnée $k$, $x_{(1),k}\leq\cdots\leq x_{(n),k}$ désigne
    les valeurs triées.

    \begin{align*}
        \texttt{Average}:\quad
        & A=\frac{1}{n}\sum_{i=1}^{n}\vect{x}_i.
        \\
        \texttt{Median}:\quad
        & M_k=\operatorname{med}_{i=1,\ldots,n}x_{i,k}.
        \\
        \texttt{TrimmedMean}:\quad
        & T_k=\frac{1}{n-2f}\sum_{r=f+1}^{n-f}x_{(r),k}.
    \end{align*}
    \texttt{TrimmedMean} retire donc les $f$ valeurs extrêmes de chaque côté,
    séparément pour chaque coordonnée.
\end{frame}

% RÉSERVE. Krum, MultiKrum et Bulyan.
\begin{frame}{Notations des agrégateurs : distances et filtrage}
    \small
    Pour chaque $i$, $\mathcal{N}_i$ contient les $n-f-2$ gradients les plus
    proches de $\vect{x}_i$ parmi les autres gradients.
    \[
        s_i=\sum_{j\in\mathcal{N}_i}\norm{\vect{x}_i-\vect{x}_j}^{2},
        \qquad i^*=\arg\min_i s_i.
    \]
    \begin{align*}
        \texttt{Krum}:\quad & \vect{x}_{i^*}.
        \\
        \texttt{MultiKrum}:\quad
        & \frac{1}{m}\sum_{i\in I_m}\vect{x}_i,
          \quad I_m=\text{indices des $m$ plus petits scores }s_i.
    \end{align*}
    Les scores sont calculés une seule fois dans la bibliothèque.

    \medskip
    \texttt{Bulyan} sélectionne itérativement $\theta=n-2f-2$ candidats avec
    \texttt{MultiKrum}. Pour chaque coordonnée $k$, il moyenne les
    $\theta-2f$ valeurs sélectionnées les plus proches de leur médiane.
    \[
        B_k=\frac{1}{\theta-2f}\sum_{r\in J_k}x_{(r),k},
        \qquad |J_k|=\theta-2f.
    \]
\end{frame}

% RÉSERVE. Méthodes géométriques et énumératives.
\begin{frame}{Notations des agrégateurs : géométrie et recherche}
    \small
    \begin{align*}
        \texttt{GeoMed}:\quad
        & \arg\min_{\vect{y}\in\mathbb{R}^d}
          \sum_{i=1}^{n}\norm{\vect{y}-\vect{x}_i}.
        \\
        \texttt{Medoid}:\quad
        & \arg\min_{\vect{x}\in\{\vect{x}_1,\ldots,\vect{x}_n\}}
          \sum_{i=1}^{n}\norm{\vect{x}-\vect{x}_i}.
        \\
        \texttt{Brute}:\quad
        & \frac{1}{|S^*|}\sum_{i\in S^*}\vect{x}_i,
          \quad S^*=\arg\min_{\substack{S\subseteq\{1,\ldots,n\}\\|S|=n-f}}
          \max_{i,j\in S}\norm{\vect{x}_i-\vect{x}_j}.
    \end{align*}
    \texttt{GeoMed} minimise la somme des distances sans imposer que la sortie
    soit un gradient reçu. \texttt{Medoid} impose cette contrainte.
    \texttt{Brute} choisit le sous-ensemble de diamètre minimal.
\end{frame}

% RÉSERVE. AKSEL et NNA.
\begin{frame}{Notations des agrégateurs : pivot et voisinage}
    \small
    \texttt{AKSEL} prend pour pivot $\vect{p}$ la médiane coordonnée par
    coordonnée, puis moyenne les $n-f$ gradients les plus proches de ce pivot.
    \[
        p_k=\operatorname{med}_{i=1,\ldots,n}x_{i,k},
        \qquad
        \operatorname{AKSEL}(\vect{x}_{1:n})=
        \frac{1}{n-f}\sum_{i\in I_{n-f}}\vect{x}_i,
        \quad I_{n-f}=\text{indices les plus proches de }\vect{p}.
    \]

    \medskip
    Pour \texttt{NNA}, $\vect{z}^{(0)}$ est le pivot et $q$ est le nombre
    de vecteurs à agréger, paramètre \texttt{num\_closest}. L'application
    $\tau$ ordonne les autres vecteurs reçus par distance croissante au pivot,
    de sorte que
    la sortie inclut le pivot et ses $q-1$ plus proches voisins :
    \[
        \operatorname{NNA}=\frac{1}{q}\sum_{r=0}^{q-1}
        \vect{z}^{(\tau(r))},
        \qquad \tau(0)=0.
    \]
    Dans \texttt{MoNNA}, le modèle local sert de pivot et $q=n-2f$.
\end{frame}

% RÉSERVE. Notations des cinq attaques implémentées.
\begin{frame}{Notations des attaques : gradients honnêtes et statistiques}
    \small
    On note $\mathcal{H}$ l'ensemble des gradients honnêtes,
    $|\mathcal{H}|=n-f$, et
    $\bar{\vect{g}}=\frac{1}{n-f}\sum_{i\in\mathcal{H}}\vect{g}_i$.
    Pour chaque coordonnée $k$, $\mu_k$ et $\sigma_k$ sont la moyenne et
    l'écart-type honnêtes.

    \begin{align*}
        \texttt{SignFlip}:\quad
        & \vect{b}=-c\,\bar{\vect{g}}.
        \\
        \texttt{ALIE}:\quad
        & b_k=\mu_k\pm z\sigma_k,
          \quad\text{avec le signe choisi pour l'attaque.}
        \\
        \texttt{Gaussian}:\quad
        & \vect{b}\sim\mathcal{N}(\mu_G\vect{1},\sigma_G^2 I_d),
          \quad\text{indépendamment des gradients honnêtes.}
    \end{align*}
    Les $f$ participants malveillants envoient chacun le vecteur ainsi défini.
    Pour \texttt{ALIE}, $z$ règle l'écart aux statistiques honnêtes. Pour
    \texttt{Gaussian}, $\mu_G$ et $\sigma_G$ sont les paramètres du bruit.
\end{frame}

% RÉSERVE. Attaques ciblées de la vulnérabilité cachée.
\begin{frame}{Notations des attaques : négation et perturbation}
    \small
    \texttt{FullGradientNegation} suppose connu le gradient calculé sur
    l'ensemble des données :
    \[
        \vect{b}=-\kappa\vect{g}_{\mathrm{full}}.
    \]

    \medskip
    \texttt{SmallPerturbation} part de la moyenne honnête et la décale dans
    une direction $\vect{e}$. L'attaque choisit le plus grand décalage qui
    reste sélectionné par l'agrégateur ciblé :
    \[
        \vect{b}=B(\gamma_m)=\bar{\vect{g}}+\gamma_m\vect{e},
        \qquad
        \gamma_m=\sup\{\gamma\geq0:
        B(\gamma)\text{ est sélectionné}\}.
    \]
    Par défaut, $\vect{e}$ est le vecteur de base de la coordonnée présentant
    la plus forte variance honnête. Les $f$ attaquants envoient ce même vecteur.
\end{frame}

% RÉSERVE. Définition formelle du modèle de menace et de la résilience.
\begin{frame}{Modèle de menace et résilience}
    \small
    Parmi les $n$ gradients reçus, au plus $f$ peuvent être arbitraires.
    Ils peuvent être coordonnés et dépendre des gradients honnêtes.
    Ces derniers sont supposés indépendants, de même loi, et d'espérance
    $\vect{g}$.

    \medskip
    Une règle $F$ est $(\alpha,f)$-résiliente si son espérance reste orientée
    vers le vrai gradient et si ses fluctuations restent contrôlées :
    \[
        \langle\mathbb{E}[F],\vect{g}\rangle
        \geq (1-\sin\alpha)\norm{\vect{g}}^2>0,
        \qquad 0\leq\alpha<\frac{\pi}{2}.
    \]
    Les moments $\mathbb{E}\norm{F}^{r}$, pour $r=2,3,4$, doivent aussi être
    bornés en fonction des moments des gradients honnêtes.

    \medskip
    Cette propriété ne suffit pas seule à prouver la convergence : le bruit
    honnête et le calendrier du pas d'apprentissage interviennent également.
\end{frame}

% RÉSERVE. Écarts des implémentations documentés dans le mémoire.
\begin{frame}{Écarts entre publication et implémentation}
    \small
    \begin{tabular}{@{}p{0.18\textwidth}p{0.39\textwidth}p{0.37\textwidth}@{}}
        \textbf{Règle} & \textbf{Publication} & \textbf{Bibliothèque} \\
        \hline
        \texttt{MultiKrum} &
        Sélection itérative, avec recalcul après chaque retrait. &
        Scores calculés une fois, puis moyenne des $m$ plus petits. \\
        \texttt{Bulyan} &
        Sélection de $n-2f$ candidats. &
        Sélection de $n-2f-2$ candidats, puis filtrage coordonnée par coordonnée.
    \end{tabular}

    \medskip
    Pour $m=1$, \texttt{MultiKrum} coïncide avec \texttt{Krum}.
    Les résultats des expériences caractérisent les variantes effectivement
    codées. Une garantie théorique ne se transfère pas automatiquement à une
    variante modifiée.
\end{frame}

% RÉSERVE. Résultats et portée des deux expériences centralisées.
\begin{frame}{Ce que montrent les deux expériences}
    \small
    \begin{block}{Spambase, paramètres de 2017}
        $n=20$, $f=6$, \texttt{MultiKrum} avec $m=14$, \texttt{SignFlip}
        d'amplitude 10. La moyenne diverge dans cette simulation ;
        \texttt{MultiKrum} poursuit l'apprentissage.
        Le réglage $m=14$ dépasse la borne théorique $m\leq5$.
    \end{block}
    \begin{block}{MNIST, attaque de 2018}
        $n=39$, $f=9$, $\gamma=4$. L'exactitude plafonne vers 90\,\% pour
        \texttt{Krum}, contre environ 94\,\% pour \texttt{Bulyan} et 95\,\%
        pour la référence saine.
    \end{block}
    Ces courbes illustrent les phénomènes visés dans ces réglages. Elles
    n'établissent ni une garantie théorique ni une stabilité statistique
    sur plusieurs graines.
\end{frame}

% RÉSERVE. Distinction entre fidélité du protocole et répétabilité.
\begin{frame}{Fidélité, répétabilité et portée}
    \small
    \begin{itemize}
        \item \textbf{Fidélité :} conserver les mécanismes utiles du protocole
        étudié et documenter les adaptations de la simulation.
        \item \textbf{Répétabilité :} une graine fixe les tirages pseudo-aléatoires.
        Dans l'environnement déterministe documenté, relancer la même
        configuration donne les mêmes résultats.
        \item \textbf{Limite :} une graine ne mesure pas la variabilité entre
        exécutions. Il faudrait plusieurs graines et une analyse statistique.
        \item \textbf{Périmètre :} les expériences tournent séquentiellement sur
        une machine ; les échanges entre workers sont simulés.
    \end{itemize}
\end{frame}
